\documentstyle[%


righttag%


,11pt%


,amssymb%


,russian%               remove in Eng.


,izvru%                 Set izven% in Eng.


]{amsart}





%





\newcommand{\la}{\langle}


\newcommand{\ra}{\rangle}


\newcommand{\argmin}{\operatornamewithlimits{arg\,min}}


\newcommand{\tts}[1]{}


\renewcommand{\baselinestretch}{1.05}





\theoremstyle{definition}


\theorembodyfont{\rm}


\newtheorem{cornonum}{\hskip\parindent Следствие}


\renewcommand{\thecornonum}{}








%\hoffset=-15mm                  For Eng.





\setcounter{page}{62}


\god{1997}


\nomer{12~(427)} %                Remove in Eng.


%\nomer{Vol.\,41, No.\,12}        Set in Eng.


% \str{75--81}                    Set in Eng.


\udk{519.853}


\author{\it И.Я.\,Заботин}


\title{Алгоритм второго порядка с параметризованными\\


направлениями для задач условной оптимизации}





\date{}


% \pagestyle{myheadings}         Set in Eng.





\pagestyle{plain} %              Remove in Eng.


\begin{document}





% \thispagestyle{plain}          Set in Eng.





\maketitle





% Body text





Предлагается алгоритм решения задачи условной минимизации сильно


выпуклой дифференцируемой функции, в котором направления спуска


задаются как линейная комбинация градиентов активных ограничений и


целевой функции. Для поиска коэффициентов комбинации решается задача


минимизации вспомогательной выпуклой квадратичной функции на


специальным образом построенном множестве. Для некоторых видов задач


математического программирования такой подход дает возможность снизить


размерность задачи построения итерационной точки по сравнению с


некоторыми известными методами условной минимизации второго порядка.





Решается задача


\begin{equation}


\min_{x\in\cal D}f_0(x),


\end{equation}


где $\cal D=\{x\in R_n:f_j(x)\le 0,\;j\in J\}$, функции


$f_j(x)$, $j\in J$, выпуклы и


непрерывно дифференцируемы в


$n$-мерном евклидовом пространстве $R_n$, функция $f_0(x)$ сильно


выпукла и


дважды непрерывно дифференцируема на $\cal D$, множество $\cal D$


удовлетворяет


условию регулярности Слейтера. Положим $J=\{1,\dotsc,m\}$,


$I=J\cup\{0\}$, $x^*=\arg\min\limits_{x\in\cal D}f_0(x)$,


$f^*_0=f_0(x^*)$, $K=\{0,1,\dotsc\}$.


Обозначим через


$f'_i(x)$, $i\in I$, и $f''_0(x)$ соответственно градиент функции


$f_i(x)$ и матрицу вторых производных функции $f_0(x)$ в точке $x$.





Пусть $\Bar x\in\cal D$, $\varepsilon\ge 0$. Положим


$J(\Bar x,\varepsilon)=\{j\in J:-\varepsilon\le f_j(\Bar x)\}$,


$I(\Bar x,\varepsilon)=\{i\in J(\Bar x,\varepsilon)\cup\{0\}:f'_i(\Bar


x)\ne 0\}$, $M(\Bar x,\varepsilon)=\Bigl\{x\in R_n:x=\Bar x+


\sum\limits_{i\in I(\Bar x,\varepsilon)}\alpha_if'_i(\Bar x),\;


\alpha_i\in R_1,\;i\in I(\Bar x,\varepsilon)\Bigr\}$,


$\cal D_1(\Bar x,\varepsilon)=\{x:x\in M(\Bar x,\varepsilon),\;


f_j(x)\le 0\;\forall j\in J\}$,


$\cal D_2(\Bar x,\varepsilon)=\{x:x\in M(\Bar x,


\varepsilon),\;


f_j(x)\le 0\;\forall j\in J(\Bar x,\varepsilon)\}$.





Отметим $\cal D_i(\Bar x,\varepsilon)\ne \varnothing$, $i=1,2$, т.к.


$\Bar x\in\cal D$. Нетрудно


проверить, что множества $\cal D_i(\Bar x,\varepsilon)$, $i=1,2$,


выпуклы и замкнуты.





Определим в $R_n$ функцию


$F(x)=\la f'_0(\Bar x),x-\Bar x\ra+\frac{1}{2}\la f''_0(\Bar x)


(x-\Bar x),x-\Bar x\ra$. Известно


(\cite{Vas}, сс.333, 184), что $F(x)$


сильно выпукла и достигает минимального значения в единственной точке


на любом выпуклом замкнутом множестве из $R_n$.





Пусть


\begin{equation}


f'_0(\Bar x)\ne 0,


\end{equation}


$\Bar y$ --- точка минимума функции $F(x)$ на множестве


$\cal D_1(\Bar x,\varepsilon)$ или $\cal D_2(\Bar x,\varepsilon)$, и


$\Bar y=\Bar x+\sum\limits_{i\in I(\Bar x,\varepsilon)}\alpha_i


(\Bar x)f'_i(\Bar x)$,


где $\alpha_i(\Bar x)\in R_1$,


$i\in I(\Bar x,\varepsilon)$. Если


$\Bar y=\argmin\limits_{x\in\cal D_1(\Bar x,\varepsilon)}F(x)$,


то числа


$\alpha_i(\Bar x)$, $i\in I(\Bar x,\varepsilon)$, очевидно,


отыскиваются


как решение задачи


\begin{gather}


\sum_{i\in I(\Bar x,\varepsilon)}\alpha_i\la f'_0(\Bar x),


f'_i(\Bar x)\ra+\frac{1}{2}\sum_{i,j\in I(\Bar x,\varepsilon)}


\alpha_i\alpha_j\la c_i(\Bar x),f'_j(\Bar x)\ra-\min,\\


f_j\Bigl(\Bar x+\sum_{i\in I(\Bar x,\varepsilon)}\alpha_i


f'_i(\Bar x)\Bigr)\le 0\quad \forall j\in J,


\end{gather}


где $c_i(\Bar x)=f''_0(\Bar x)f'_i(\Bar x)$,


$i\in I(\Bar x,\varepsilon)$. Пусть


$\Bar y=\argmin\limits_{x\in\cal D_2(\Bar x,\varepsilon)}F(x)$. Если


$J(\Bar x,\varepsilon)\ne\varnothing$, то числа


$\alpha_i(\Bar x)$, $i\in I(\Bar x,\varepsilon)$, находятся как


решение задачи (3) при ограничениях


\begin{equation}


f_j\Bigl(\Bar x+\sum_{i\in I(\Bar x,\varepsilon)}\alpha_if'_i(\Bar x)


\Bigr)\le 0\quad \forall j\in J(\Bar x,\varepsilon).


\end{equation}


Если $J(\Bar x,\varepsilon)=\varnothing$, то $\alpha_0(\Bar x)$ --- точка


минимума одномерной функции


$\alpha_0\la f'_0(\Bar x),f'_0(\Bar x)\ra+\linebreak


\frac{1}{2}\alpha^2_0 \la c_0(\Bar x),f'_0(\Bar x)\ra$


на прямой $\cal D_2(\Bar x,\varepsilon)=M(\Bar x,\varepsilon)=


\{x\in R_n:x=\Bar x+\alpha_0f'_0(\Bar x),\;\alpha_0\in R_1\}$.





Предположим, что определенная выше точка $\Bar y$ удовлетворяет


соотношениям


$F(\Bar y)<F(\Bar x)=0$. Покажем, что тогда $s=\Bar y-\Bar x$ ---


подходящее для функции $f_0(x)$ в точке $\Bar x$


направление. Поскольку для сильно выпуклой функции $f(x)$ выполняется


(\cite{Vas}, с.185) неравенство


\begin{equation}


\la f''_0(x)\xi,\xi\ra\ge \mu\|\xi\|^2\quad\forall x\in\cal D,


\quad \xi\in R_n,


\end{equation}


где $\mu>0$, то $\la f''_0(\Bar x)s,s\ra >0$, а значит,


$\la f'_0(\Bar x),s\ra<0$, и $s$ ---


направление спуска для функции $f_0(x)$


в точке $\Bar x$. Осталось показать, что $s$ --- возможное в точке


$\Bar x$


направление. Будем считать


\begin{equation}


J(\Bar x,0)\ne\varnothing,


\end{equation}


т.к. в противном случае утверждение очевидно. Тогда согласно (4), (5)


$f_j(\Bar x+s)\le 0$ $\forall j\in J(\Bar x,0)$. Поскольку


$f_j(\Bar x)<0$ $\forall j\in J\setminus J(\Bar x,0)$, то


найдется такое число $\gamma\in(0,1]$, что


$f_j(\Bar x+\gamma s)\le 0$ $\forall j\in J$, и


утверждение доказано. Такой способ выбора подходящих направлений будет


использован ниже в предлагаемом методе решения задачи (1).





Заметим, что в задачах (3), (4) и (3), (5) выбора направления в точке


$\Bar x$ число переменных зависит лишь от


$\mu(I(\Bar x,\varepsilon))$ ---


количества индексов в


множестве $I(\Bar x,\varepsilon)$, а не от размерности пространства


$R_n$.


Поэтому при $\mu(I(\Bar x,\varepsilon))<n$


эти задачи имеют меньше переменных, чем, например, аналогичная задача


построения в точке $\Bar x$ подходящего направления вида


$s=\arg\min\limits_{x\in\cal D}F(x)-\Bar x$,


используемая в


методе Ньютона \cite{Vas}. Кроме того, задача (3), (5) имеет и меньшее


число ограничений по сравнению с задачей выбора направления в


упомянутом известном методе.





Непосредственно из определения строгой выпуклости функции $f(x)$


следует





\begin{lem}


Пусть $f(x)$ --- строго выпуклая в $R_n$ функция,


$\Bar x\in R_n$, $a_i\in R_n$, $a_i\ne 0$ $\forall i=1,\dotsc,p$.


Тогда


функция $\varphi(\alpha)=f\Bigl(\Bar x+\sum\limits^p_{i=1}


\alpha_i a_i\Bigr)$, где $\alpha=(\alpha_1,\dotsc,\alpha_p)$,


$\alpha_i\in R_1$


$\forall i=1,\dotsc,p$, строго выпукла по переменным $\alpha$.


\end{lem}





\begin{thm}


Пусть $\Bar x\in\cal D$, $\varepsilon\ge 0$, выполняется условие {\em


(2)} и $\Bar y$ --- точка минимума


функции $F(x)$ на множестве $\cal D_1(\Bar x,\varepsilon)$ или


$\cal D_2(\Bar x,\varepsilon)$. Для


того чтобы точка $\Bar x$


была решением задачи {\em (1)}, необходимо и достаточно равенства


\begin{equation}


F(\Bar y)=0.


\end{equation}


\end{thm}





\begin{pf*}{Доказательство. Необходимость}


Пусть $\Bar x$ --- точка минимума функции $f_0(x)$ на множестве


$\cal D$. Так


как $F(\Bar x)=0$, $\Bar x\in\cal D_i(\Bar x,\varepsilon)$, $i=1,2$,


то $F(\Bar y)\le 0$. Допустим, что последнее неравенство выполняется


как строгое.


Тогда, как показано выше, направление $\Bar y-\Bar x$


является подходящим


для функции $f_0(x)$ в точке $\Bar x$, что противоречит выбору


$\Bar x$.


Следовательно, утверждение (8) доказано.


\medskip





{\bf Достаточность.} Пусть выполняется равенство (8). Положим


$f'_i(\Bar x)=a_i$, $c_i(\Bar x)=c_i$,


$i\in I$, и будем считать, что


$\Bar y=\Bar x+\sum\limits_{i\in I(\Bar x,\varepsilon)}\alpha^*_ia_i$,


где $\alpha^*_i\in R_1$ $\forall i\in I(\Bar x,\varepsilon)$. Докажем


сначала, что справедливо


(7). Допустим противное, т.е.


$f_j(\Bar x)<0$, $j\in J$. Поскольку функции $f_j(x)$ непрерывны и


$\la c_0,a_0\ra>0$ согласно (6), то найдется такое достаточно близкое к


нулю отрицательное число $\Bar \alpha_0$, что


\begin{equation}


f_j(\Bar x+\Bar \alpha_0a_0)\le 0\quad \forall j\in J,\quad


\Bar \alpha_0\la a_0,a_0\ra+\frac{1}{2}\Bar \alpha^2_0\la c_0,a_0\ra<0.


\end{equation}


Положим $\Bar \alpha_i=0$ $\forall i\in I(\Bar x,\varepsilon)\setminus


\{0\}$,


$\Bar{\Bar x}=\Bar x+\sum\limits_{i\in I(\Bar x,\varepsilon)}\Bar


\alpha_ia_i$. Тогда в силу первого из неравенств (9) выполняется


включение $\Bar{\Bar x}\in\cal D_i(\Bar x,\varepsilon)$, $i=1,2$.


Значит, по условию теоремы


$F(\Bar y)\le F(\Bar{\Bar x})=\sum\limits_{i\in I(\Bar x,\varepsilon)}


\Bar \alpha_i\la a_0,a_i\ra+\frac{1}{2}\sum\limits_{i,j\in I(\Bar


x,\varepsilon)}\Bar \alpha_i\Bar\alpha_j\la c_i,a_j\ra=\Bar\alpha_0


\la a_0,a_0\ra+\frac{1}{2}\Bar\alpha^2_0\la c_0,a_0\ra$.


Отсюда с учетом


второго из соотношений (9) следует неравенство $F(\Bar y)<0$,


противоречащее условию (8). Таким образом, утверждение (7) доказано.





Далее, множество $I(\Bar x,\varepsilon)$ кроме индекса $i=0$ содержит


согласно условию


Слейтера также все индексы $i\in J(\Bar x,0)$. Не ограничивая общности,


будем считать,


что $J(\Bar x,0)=\{1,\dotsc,r\}$,


$I(\Bar x,\varepsilon)=\{0,1,\dotsc,l\}$, где


$1\le r\le l\le m$. Положим


$\alpha=(\alpha_0,\dotsc,\alpha_l)$, $\alpha_i\in R_1$,


$i=0,\dotsc,l$, $\Lambda_1=\Bigl\{\alpha:f_j\Bigl(\Bar x+


\sum\limits^l_{i=0}\alpha_ia_i\Bigr)\le 0\;\forall j\in J\Bigr\}$


--- допустимое множество


задачи (3), (4),


$\Lambda_2=\Bigl\{\alpha:f_j\Bigl(\Bar x+


\sum\limits^l_{i=0}\alpha_ia_i\Bigr)\le 0\;\forall j\in


J(\Bar x,\varepsilon)\Bigr\}$


--- допустимое множество задачи (3), (5). Так как


$0\in\Lambda_i$, $i=1,2$, а по условию теоремы


$\alpha^*=(\alpha^*_0,\dotsc,\alpha^*_l)$ --- решение задачи (3),


(4) или (3),


(5), причем в силу леммы 1 это решение единственно, то согласно (8)


$\alpha^*=0$.


Нетрудно проверить, что в точке $\alpha=0$ градиентами целевой функции


(3) и функций ограничений


$f_j\Bigl(\Bar x+\sum\limits^l_{i=0}\alpha_ia_i\Bigr)$,


$j=1,\dotsc,r$, задачи (3),


(4) и (3), (5) являются


соответственно векторы


$d_0=(\la a_0,a_0\ra,\dotsc,\la a_0,a_l\ra)$,


$d_j=(\la a_j,a_0\ra,\dotsc,\la a_j,a_l\ra)$, $j=1,\dotsc,r$. Поскольку


$f_j\Bigl(\Bar x+\sum\limits^l_{i=0}\alpha^*_ia_i\Bigr)=


f_j(\Bar x)=0$, $j=1,\dotsc,r$, и по


лемме 1


(\cite{Zab}, с.18) множества $\Lambda_1$, $\Lambda_2$ удовлетворяют


условию Слейтера, то


найдутся такие числа $\gamma_j\le 0$, $j=1,\dotsc,r$, что


$d_0=\sum\limits^r_{j=1}\gamma_jd_j$ (\cite{Kar},


с.50). Отсюда следуют равенства


\begin{equation}


\la a_0,a_j\ra=\la b,a_j\ra,\quad j=0,\dotsc,l,


\end{equation}


где


\begin{equation}


b=\sum^r_{j=1}\gamma_ja_j.


\end{equation}


Умножив равенства (10) на $\gamma_j$ и просуммировав их по $j$ от 1 до


$l$,


получим $\la a_0,b\ra=\la b,b\ra$. Отсюда с учетом (10) следует, что


$\la a_0,b\ra^2=\la a_0,a_0\ra\la b,b\ra$,


поэтому векторы $a_0$ и


$b$ линейно зависимы, т.е. $a_0=\gamma b$. Так как из (10) и (2)


следует неравенство


$\la a_0,b\ra >0$, то $\gamma>0$. Тогда в силу (11)


$f'_0(\Bar x)=\sum\limits^r_{j=1}\gamma\gamma_jf'_j(\Bar x)$. Положив


$\tau_j=\gamma\gamma_j$, $j=1,\dotsc,r$; $\tau_j=0$,


$j=r+1,\dotsc,m$, получим отсюда


равенства


$f'_0(\Bar x)=\sum\limits^m_{j=1}\tau_jf'_j(\Bar x)$,


$\sum\limits^m_{j=1}\tau_jf_j(\Bar x)=0$.


Следовательно, по теореме Куна-Таккера


(\cite{Kar}, с.54) $\Bar x$ --- точка минимума функции $f_0(x)$ на


множестве $\cal D$.


\end{pf*}





\begin{cornonum}


Пусть выполняются условия теоремы 1. Для того чтобы точка $\Bar x$ была


решением задачи (1), необходимо и достаточно равенства


\begin{equation}


\Bar y=\Bar x.


\end{equation}


\end{cornonum}





\begin{pf}


Пусть $\Bar x$ --- решение задачи (1). Допустим, что (12) не


выполняется. По


теореме 1 справедливо (8), а значит, $F(\Bar x)=F(\Bar y)$.


Тогда в силу выпуклости


множеств $\cal D_i(\Bar x,\varepsilon)$, $i=1,2$, и сильной выпуклости


функции $F(x)$ направление $\Bar x-\Bar y$ является


подходящим для функции $F(x)$ в точке $\Bar y$,


что противоречит выбору


точки $\Bar y$. Следовательно, равенство (12) доказано.





Пусть справедливо (12). Поскольку $F(\Bar x)=0$, то выполняется (8), а


значит, по теореме 1 точка $\Bar x$ --- решение задачи (1).


\end{pf}





Предлагаемый алгоритм решения задачи (1) заключается в следующем.


Выбирается точка начального приближения $x_0\in\cal D$ и числа


$\Bar\varepsilon>0$, $1<\theta<+\infty$.


Пусть уже построена точка $x_k$, $k\in K$.


\begin{itemize}


\item[1.] Если $f'_0(x_k)=0$, то $x_k$ --- решение задачи, и процесс


заканчивается.


\item[2.] При $\varepsilon_k\ge\Bar\varepsilon$ отыскивается $y_k$ ---


точка минимума функции


$$


F_k(x)=\la f'_0(x_k),x-x_k\ra+\frac{1}{2}\la


f''_0(x_k)(x-x_k),x-x_k\ra


$$


на множестве $\cal D_1(x_k,\varepsilon_k)$ или


$\cal D_2(x_k,\varepsilon_k)$. Если $F_k(y_k)=0$, то $x_k=x^*$, и


процесс заканчивается.


\item[3.] Полагается $p_k=y_k-x_k$, $s_k=\tau_kp_k$, где $\tau_k=1$,


если


$y_k\in\cal D$, и $\tau_k$ выбирается из условий


\begin{equation}


\tau_k>0,\quad x_k+\tau_kp_k\in\cal D;\quad


x_k+\theta\tau_kp_k\notin\cal D\text{\; в противном случае.}


\end{equation}


\item[4.] Полагается $x_{k+1}=x_k+\beta_ks_k$. Числа $\beta_k$


выбираются


либо из условия


(\cite{Ps}, с.209; \cite{Vas}, с.331) $\beta_k=\lambda^{i_k}$


$\forall k\in K$, где $i_k$


--- наименьший из номеров $i\ge 0$, при которых выполняется


неравенство


\begin{equation}


f_0(x_k)-f_0(x_k+\lambda^is_k)\ge \eta\lambda^i\tau_k|F_k(y_k)|,


\end{equation}


а параметры $\lambda$, $\eta$ таковы, что $0<\lambda<1$,


$0<\eta<1$, либо из условий


(\cite{Kar}, с.212)


\begin{equation}


\begin{split}


0\le\beta\le 1,\quad &f_0(x_k+\beta_ks_k)\le(1-\nu_k)f_0(x_k)+


\nu_k\omega_k,\\


0<\nu\le\nu_k\le 1,\quad &\omega_k=\min_{0\le\beta\le 1}


f_0(x_k+\beta s_x)\quad \forall k\in K.


\end{split}


\end{equation}


\end{itemize}





Отметим, что в алгоритме $I(x_k,\varepsilon_k)\ne\varnothing$


$\forall k\in K$, т.к. $f'_0(x_k)\ne 0$


$\forall k\in K$. Если $x_k$ --- решение задачи


(1), то выход из цикла в п.2 происходит в соответствии с теоремой 1, в


противном случае, как показано выше, $p_k$ --- подходящее для


функции $f_0(x)$ в точке $x_k$ направление, и выбор чисел


$\tau_k$ возможен.


Если


точка $y_k$ для всех $k\in K$ отыскивается с помощью решения задачи (3),


(4)


при $\Bar x=x_k$, $\varepsilon=\varepsilon_k$, то $y_k\in\cal D$


$\forall k\in K$, и


задача нахождения чисел $\tau_k$ из условий (13) не возникает.





Возможность отыскания шагов $\beta_k$ из условий (14) будет обоснована


позже.


Подчеркнем также, что числа $\beta_k$ выбираются в алгоритме для всех


$k\in K$


только одним из предложенных способов, т.е. чередование способов в


процессе построения последовательности $\{x_k\}$


не предполагается. Позже при


дополнительных предположениях будет обсуждена возможность использования


в алгоритме и других способов задания шагов.





\begin{lem}


Последовательности $\{x_k\}$ и $\{p_k\}$, построенные по предложенному


алгоритму, ограничены.


\end{lem}





\begin{pf}


Так как функция $f_0(x)$ сильно выпукла и непрерывна на $\cal D$, то


множество $M(x_0)=\{x\in\cal D:f_0(x)\le f_0(x_0)\}$


ограничено. Поскольку $y_k\ne x_k$ $\forall k\in K$, то, как показано


выше, направление $p_k$ для


всех $k\in K$ является подходящим для функции $f_0(x)$ в точке $x_k$.


Тогда согласно (14), (15) $x_k\in M(x_0)$ $\forall k\in K$, и


последовательность $\{x_k\}$ ограничена.





Далее, из последнего утверждения следует


существование такой константы $d>0$, что


\begin{equation}


\|f'_0(x_k)\|\le d\quad \forall k\in K.


\end{equation}


Поскольку


$F_k(y_k)=\la f'_0(x_k),p_k\ra+\frac{1}{2}\la f''_0(x_k)p_k,p_k\ra<


F_k(x_k)=0$


$\;\forall k\in K$, то отсюда с учетом (16), (6) следует неравенство


$-d\|p_k\|+\frac{\mu}{2}\|p_k\|^2<0$ $\forall k\in K$.


Тогда $\|p_k\|<2d/\mu$ $\;\forall k\in K$.


\end{pf}





\begin{lem}


Пусть $\Bar x$ --- предельная точка последовательности $\{x_k\}$,


построенной по


предложенному алгоритму, и $x_k\to\Bar x$ при $k\to\infty$,


$k\in K'\subset K$. Тогда


найдется такой номер $N$, что для всех $k\ge N$, $k\in K'$,


справедливо включение


\begin{equation}


I(\Bar x,0)\subset I(x_k,\varepsilon_k).


\end{equation}


\end{lem}





\begin{pf}


Если $J(\Bar x,0)=\varnothing$, то $I(\Bar x,0)=\{0\}$ либо


$I(\Bar x,0)=\varnothing$, а


поскольку по алгоритму


\begin{equation}


0\in I(x_k,\varepsilon_k)\quad \forall k\in K',


\end{equation}


то включение (17) справедливо при всех $k\in K'$. Поэтому далее будем


считать,


что выполняется (7). Тогда по условию регулярности Слейтера


$f'_i(\Bar x)\ne 0$ $\forall i\in J(\Bar x,0)$, и,


следовательно, $I(\Bar x,0)=\{0\}\cup J(\Bar x,0)$, либо


$I(\Bar x,0)=J(\Bar x,0)$. Согласно (18)


остается показать включение


$J(\Bar x,0)\subset I(x_k,\varepsilon_k)$ для всех достаточно больших


номеров $k\in K'$. Пусть $l\in J(\Bar x,0)$,


$E=\{x\in\cal D:-\Bar\varepsilon\le f_l(x)\}$. Поскольку


$f'_l(\Bar x)\ne 0$ и


функция $f_l(x)$ непрерывно дифференцируема, то найдется такая


окрестность


$\omega$ точки $\Bar x$, что


$f'_l(x)\ne 0$ $\forall x\in\omega\cap E$. Выберем номер $N$


так, чтобы $x_k\in\omega\cap E$


$\forall k\ge N$, $k\in K'$. Тогда


$-\varepsilon_k\le\Bar\varepsilon\le f_l(x_k)\le 0$,


$k\ge N$, $k\in K'$, а значит,


$l\in J(x_k,\varepsilon_k)$,


$k\ge N$, $k\in K'$. Поскольку


$f'_l(x_k)\ne 0$, $k\ge N$, $k\in K'$, то


$l\in I(x_k,\varepsilon_k)$, $k\ge N$, $k\in K'$.


\end{pf}





\begin{lem}


Для последовательности $\{\tau_k\}$, $k\in K$, определенной согласно


п.{\em 3} алгоритма, существует такое число $\tau'>0$, что


$\tau_k\ge \tau'$ $\forall k\in k$.


\end{lem}





\begin{pf}


Допустим, что утверждение леммы неверно, т.е. существует такая


подпоследовательность $\{\tau_k\}$, $k\in K_1\subset K$, что


\begin{equation}


\tau_k\to 0,\quad k\to\infty,\quad k\in K_1.


\end{equation}


Поскольку $\tau_k\in(0,1]$ $\forall k\in K$, и $\tau_k=1$


лишь в случае,


когда $y_k\in\cal D$, то в силу предположения (19) можно считать, что


\begin{equation}


y_k\notin\cal D\quad \forall k\in K_1.


\end{equation}


Тогда при всех $k\in K_1$ справедливы соотношения (13), а значит, для


произвольно фиксированного номера $k\in K_1$ найдется такой индекс


$i\in J$, что


\begin{equation}


f_i(x_k+\tau_kp_k)\le 0,\quad


f_i(x_k+\theta\tau_kp_k)>0,\quad f_i(x_k+p_k)>0.


\end{equation}


Так как множество $J$ конечно, то отсюда следует существование индекса


$i_0\in J$ и подпоследовательности $K_2\subset K_1$ таких, что для


всех $k\in K_2$ выполняются


неравенства (21) при $i=i_0$. Согласно лемме 2 из последовательности


$\{x_k\}$,


$k\in K_2$, можно выделить сходящуюся подпоследовательность


$\{x_k\}$, $k\in K_3\subset K_2$.


Пусть $\Bar x$


--- ее предельная точка. Тогда, переходя к пределу в первых двух


неравенствах (21) при $i=i_0$, $k\to\infty$, $k\in K_3$,


с учетом (19) и


леммы 2 получим


равенство $f_{i_0}(\Bar x)=0$. Следовательно,


$i_0\in J(\Bar x,0)$, а поскольку $f'_{i_0}(\Bar x)\ne 0$ в силу


регулярности


$\cal D$, то $i_0\in I(\Bar x,0)$. Значит, по лемме 3 найдется такой


номер $N$, что


\begin{equation}


i_0\in J(x_k,\varepsilon_k)\quad \forall k\ge N,\quad k\in K_3.


\end{equation}


Далее, согласно (20) $y_k$ --- точка минимума $F_k(x)$


на множестве $\cal D_2(x_k,\varepsilon_k)$ для


всех


$k\ge N$, $k\in K_3$, т.к. при втором способе задания $y_k$


выполняется включение $y_k\in\cal D$.


Таким образом, $y_k=x_k+p_k\in\cal D_2(x_k,\varepsilon_k)$,


$k\ge N$, $k\in K_3$. Тогда из


(22)


следует неравенство


$f_{i_0}(x_k+p_k)\le 0$, $k\ge N$, $k\in K_3$,


которое противоречит третьему из соотношений (21) при


$i=i_0$, $k\ge N$, $k\in K_3$.


Значит, предположение (19)


неверно.


\end{pf}





\begin{lem}


Пусть для множества $\cal D$ и функций $f_i(x)$, $i\in I$, выполняются


указанные выше условия, и, кроме того,


\begin{equation}


\la f''_0(x)\xi,\xi\ra\le M\|\xi\|^2\quad


\forall x\in\cal D,\quad \xi\in R_n,


\end{equation}


где $M\ge \mu$, константа $\mu$ определена в {\em (6)}. Если


последовательность $\{x_k\}$


построена по предложенному алгоритму, то существует конечный предел


\begin{equation}


\lim_{k\to\infty}f_0(x_k)\ge f^*_0,


\end{equation}


и выполняется равенство


\begin{equation}


\lim_{k\to\infty}\|p_k\|=0.


\end{equation}


\end{lem}





\begin{pf}


Если $y_k=x_k$ для некоторого $k\in K$, то по следствию к теореме 1


$x_k$ ---


решение


задачи (1). Поэтому будем считать, что $y_k\ne x_k$, а значит,


$s_k\ne 0$ и $F_k(y_k)<0$ для


всех $k\in K$.





Покажем сначала, что


\begin{equation}


f_0(x_k)-f_0(x_k+\beta s_k)\ge \beta(1-\beta M/\mu)\tau_k


|F_k(y_k)|


\end{equation}


$\forall k\in K$, $\beta\in[0,1]$. Так как по алгоритму


$\tau_k\in(0,1]$, $k\in K$, то из


выпуклости функции $F_k(x)$ при произвольном $\beta\in[0,1]$ следуют


соотношения


\begin{equation}


F_k(x_k+\beta s_k)=F_k(x_k+\beta\tau_k(y_k-x_k))\le


\beta\tau_kF_k(y_k)+(1-\beta\tau_k)F_k(x_k)=


\beta\tau_kF_k(y_k),\quad k\in K.


\end{equation}


Отсюда с учетом равенства


\begin{equation}


f_0(x_k+\beta s_k)-f_0(x_k)=F_k(x_k+\beta s_k)+


(\beta^2/2)\la[f''_0(x_k+\theta'\beta s_k)-


f''_0(x_k)]s_k,s_k\ra,


\end{equation}


где $0\le\theta'\le 1$, и неравенств (6), (23) получим оценку


\begin{equation}


f_0(x_k+\beta s_k)-f_0(x_k)\le\beta\tau_kF_k(y_k)+


(\beta^2/2)M\|s_k\|^2,


\end{equation}


$k\in K$, $\beta\in[0,1]$. Поскольку $y_k$ --- точка минимума сильно


выпуклой функции $F_k$ на


множестве $\cal D_1(x_k,\varepsilon_k)$ или


$\cal D_2(x_k,\varepsilon_k)$, а


$x_k\in\cal D_i(x_k,\varepsilon_k)$, $i=1,2$, то (\cite{Vas},


сс.182, 185)


$\|x_k-y_k\|^2\le(2/\mu)(F_k(x_k)-F_k(y_k))=(2/\mu)|F_k(y_k)|$,


и, значит,


\begin{equation}


\|s_k\|^2=\tau^2_k\|p_k\|^2\le(2/\mu)\tau_k|F_k(y_k)|,


\end{equation}


$k\in K$. Подставляя эту оценку в (29), получим (26).





Докажем теперь утверждения леммы отдельно для каждого из предложенных


способов выбора шагов $\beta_k$. Пусть при построении


последовательности $\{x_k\}$


шаги выбирались из условия (14). Покажем прежде всего, что при


произвольно фиксированном $k\in K$ найдется хотя бы один номер


$i\ge 0$, для


которого справедливо (14). Действительно, при всех значениях $\beta$,


удовлетворяющих условиям


\begin{equation}


0<\eta_0=\lambda(1-\eta)\mu/M\le\beta\le(1-\eta)\mu/M<1,


\end{equation}


согласно (26) выполняется неравенство


\begin{equation}


f_0(x_k)-f_0(x_k+\beta s_k)\ge \beta\eta\tau_k|F_k(y_k)|.


\end{equation}


Выберем номер $r_k\ge 1$ так, чтобы


$\lambda^{r_k}\le(1-\eta)\mu/M<\lambda^{r_k-1}$. Тогда


\begin{equation}


\eta_0<\lambda^{r_k}\le(1-\eta)\mu/M,


\end{equation}


и число $\beta=\lambda^{r_k}$ удовлетворяет условиям (31), а значит, и


неравенству


(32). Следовательно, (14) выполняется при $i=r_k$, и среди номеров


$i\le r_k$,


удовлетворяющих (14), можно выбрать наименьший номер $i=i_k$. Таким


образом,


возможность выбора шагов из условия (14) доказана. Далее, поскольку


$f_0(x_k)\ge f_0(x_{k+1})\ge f^*_0$,


$k\in K$, то утверждение (24) леммы справедливо, и, значит,


\begin{equation}


\lim_{k\to\infty}(f_0(x_k)-f_0(x_{k+1}))=0.


\end{equation}


Так как по алгоритму


$f_0(x_k)-f_0(x_{k+1})\ge\eta\beta_k\tau_k|F_k(y_k)|$ и согласно (33)


$\beta_k=\lambda^{i_k}\ge\lambda^{r_k}>\eta_0$ для


всех $k\in K$, то


$$


f_0(x_k)-f_0(x_{k+1})\ge\eta\eta_0\tau_k|F_k(y_k)|\quad \forall k\in K.


$$


Отсюда с учетом (34), (30) и леммы 4 следует (25).





Пусть теперь последовательность $\{x_k\}$ построена с выбором шагов из


условия


(15). Используя (26), для любого $\beta\in[0,1]$ получим неравенства


$f_0(x_k)-f_0(x_{k+1})\ge\nu_k(f_0(x_k)-f_0(x_k+\beta s_k))\ge


\nu_k\tau_k(\beta-\beta^2M/\mu)|F_k(y_k)|$,


$k\in K$. Отсюда


при $\beta=\mu/2M<1$ следует оценка


$f_0(x_k)-f_0(x_{k+1})\ge(\mu/4M)\nu\tau_k|F_k(y_k)|$ $\forall k\in K$.


Так как последовательность $\{f_0(x_k)\}$ монотонно


убывает и ограничена снизу, то справедливо утверждение (24) леммы и


равенство (34). Значит, из последнего неравенства и оценки (30) с


учетом леммы 4 следует (25).


\end{pf}





\begin{lem}


Пусть $H$ --- некоторое множество индексов, $0<\sigma'\le 1$,


$\sigma''>0$, $V$ ---


выпуклая


оболочка векторов $a_i\in R_n$, $i\in H$, и $\Bbb O\notin V$, где


$\Bbb O$ ---


нуль пространства $R_n$. Тогда


для каждого непустого подмножества $H_1\subset H$ найдутся такие числа


$\alpha_i\le 0$, $i\in H_1$, что $\sum\limits_{i\in H_1}\alpha_i\ge -1$


и


$$


\sigma'\Bigl\la a_j,\sum_{i\in H_1}\alpha_ia_i\Bigr\ra+\sigma''


\Bigl\|\,\sum_{i\in H_1}\alpha_ia_i\Bigr\|^2<0


\quad \forall j\in H_1.


$$


\end{lem}





\begin{pf}


Пусть $V_1$ --- выпуклая оболочка векторов $a_i$, $i\in H_1$, $q_1$


---


проекция $\Bbb O$ на


множество $V_1$. Поскольку $V_1\subset V$, то $\Bbb O\notin V_1$ в силу


условия леммы, и $\|q_1\|>0$. Так как


$q_1\in V_1$, то найдутся такие числа $\gamma_i\ge 0$, что


$q_1=\sum\limits_{i\in H_1}\gamma_ia_i$ и


$\sum\limits_{i\in H_1}\gamma_i=1$. По


известному свойству


проекции $\la q_1,x-q_1\ra\ge 0$ $\forall x\in V_1$. Отсюда


$\sigma'\la\sigma' q_1,x\ra\ge\la\sigma' q_1,\sigma' q_1\ra$


$\forall x\in V_1$. Положим


$\sigma' q_1=q'_1$ и выберем число $\gamma\in(0,1]$ так,


чтобы $\gamma^2\sigma''<\gamma$. Тогда для всех $x\in V_1$,


в том числе и


для всех векторов $a_j$, $j\in H_1$, выполняются соотношения


$$


\sigma'\gamma\la q'_1,x\ra\ge\gamma\la q'_1,q'_1\ra >


\gamma^2\sigma''\|q'_1\|^2=\sigma''\|\gamma q'_1\|^2.


$$


Следовательно, для всех $j\in H_1$ справедливо неравенство


$\sigma'\la a_j,\gamma q'_1\ra>\sigma''\|\gamma q'_1\|^2$, или


$$


\sigma'\Bigl\la a_j,\sum_{i\in H_1}\sigma'\gamma\gamma_ia_i


\Bigr\ra >\sigma''\Bigl\|\,\sum_{i\in H_1}(-\sigma'\gamma\gamma_i)


a_i\Bigr\|^2.


$$


Положив $-\sigma'\gamma\gamma_i=\alpha_i$, $i\in H_1$, получим отсюда


утверждение леммы.


\end{pf}





\begin{thm}


Пусть выполняются условия леммы {\em 5}. Тогда последовательность


$\{x_k\}$,


построенная по предложенному алгоритму, сходится к точке $x^*$.


\end{thm}





\begin{pf}


Пусть числа $\alpha_i(x_k)$, $i\in I(x_k,\varepsilon_k)$, таковы, что


$y_k=x_k+\sum\limits_{i\in I(x_k,\varepsilon_k)}\alpha_i(x_k)


f'_i(x_k)$, $k\in K$.


Тогда из (25) с учетом (16) следуют предельные соотношения


\begin{equation}


\lim_{k\to\infty}\Bigl\|\sum_{i\in I(x_k,\varepsilon_k)}


\alpha_i(x_k)f'_i(x_k)\Bigr\|=0,\quad


\lim_{k\to\infty}\,\sum_{i\in I(x_k,\varepsilon_k)}


\la f'_0(x_k),f'_i(x_k)\ra


\alpha_i(x_k)=0.


\end{equation}





Согласно лемме 5 для последовательности $\{x_k\}$ справедливо (24).


Допустим,


что неравенство (24) выполняется как строгое, т.е. утверждение теоремы


неверно. Выберем из ограниченной последовательности $\{x_k\}$,


$k\in K$,


сходящуюся


подпоследовательность $\{x_k\}$, $k\in K_1\subset K$, и пусть $\Bar x$


--- ее предельная точка. По


предположению $\Bar x\ne x^*$. Значит, выполняется (2),


$I(\Bar x,0)\ne\varnothing$, и существует


такой


вектор $s\in R_n$, что $\la f'_i(\Bar x),s\ra<0$


$\forall i\in I(\Bar x,0)$. Тогда нуль


пространства не принадлежит выпуклой


оболочке векторов $f'_i(\Bar x)$, $i\in I(\Bar x,0)$, и по лемме 6


найдутся такие числа $\alpha'_i\le 0$, $i\in I(\Bar x,0)$, что


\begin{equation}


\Bigl\la f'_j(\Bar x),\,\sum_{i\in I(\Bar x,0)}\alpha'_if'_i(\Bar x)


\Bigr\ra +\frac{M}{2}\Bigl\|\,


\sum_{i\in I(\Bar x,0)}\alpha'_if'_i(\Bar x)\Bigr\|^2<0\quad


\forall j\in I(\Bar x,0).


\end{equation}


Положим $p=\sum\limits_{i\in I(\Bar x,0)}\alpha'_if'_i(\Bar x)$.


Согласно


(36)


$p\ne 0$. Покажем, что при некотором значении $0<\gamma'<1$


справедливы неравенства


\begin{equation}


f_j(\Bar x+\gamma'p)<0\quad \forall j\in J.


\end{equation}


Если $J(\Bar x,0)=\varnothing$, то $f_j(\Bar x)<0$ $\forall j\in J$, и


существование числа $\gamma'$, удовлетворяющего (37),


очевидно. Пусть выполняется (7). В силу регулярности допустимого


множества $f'_j(\Bar x)\ne 0$ $\forall j\in J(\Bar x,0)$, т.е.


$J(\Bar x,0)\subset I(\Bar x,0)$.


Тогда согласно (36) $p$ ---


направление спуска


для функций $f_j(x)$, $j\in J(\Bar x,0)$, в точке $\Bar x$. Значит,


существует такое число $\Bar\gamma\in (0,1)$, что


$f_j(\Bar x+\gamma p)<f_j(\Bar x)=0$ $\forall j\in J(\Bar x,0)$,


$\gamma\in(0,\Bar\gamma]$, а


поскольку $f_j(\Bar x)<0$ $\forall j\in J\setminus J(\Bar x,0)$, то


при некотором $\gamma'\in(0,\Bar\gamma]$ выполняется (37).


Таким образом, справедливость утверждения (37) доказана при условии


(7). Далее, с учетом выбора $\gamma'$ из (36) при $j=0$ следует


неравенство


$\gamma'\la f'_0(\Bar x),p\ra+\frac{M}{2}(\gamma')^2\|p\|^2<0$. Тогда,


положив $\gamma'\alpha'_i=\Bar\alpha_i$, $i\in I(\Bar x,0)$, получим


отсюда и из (37)


\begin{gather*}


\Bigl\la f'_0(\Bar x),\,\sum_{i\in I(\Bar x,0)}\Bar\alpha_i


f'_i(\Bar x)\Bigr\ra+\frac{M}{2}\Bigl\|\,\sum_{i\in I(\Bar x,0)}


\Bar\alpha_if'_i(\Bar x)\Bigr\|^2=\delta<0,\\


f_j\Bigl(\Bar x+\sum_{i\in I(\Bar x,0)}\Bar\alpha_if'_i(\Bar x)\Bigr)


<0\quad\forall j\in J.


\end{gather*}


Так как функции $f_j(x)$, $j\in J$, непрерывно дифференцируемы на


$\cal D$,


множество $\cal D$


регулярно и по предположению выполняется (2), то найдется такая


окрестность $\omega$ точки $\Bar x$, что для всех


$x\in\omega\cap\cal D$ справедливы


соотношения


\begin{gather*}


\Bigl\la f'_0(x),\,\sum_{i\in I(\Bar x,0)}\Bar\alpha_if'_i(x)


\Bigr\ra+\frac{M}{2}\Bigl\|\,\sum_{i\in I(\Bar x,0)}


\Bar\alpha_if'_i(x)\Bigr\|^2\le \frac{\delta}{2},\\


f_j\Bigl(x+\sum_{i\in I(\Bar x,0)}\Bar\alpha_if'_i(x)\Bigr)\le 0


\quad \forall j\in J,\qquad f'_i(x)\ne 0\quad \forall i\in


I(\Bar x,0).


\end{gather*}


Поскольку $x_k\to\Bar x$, $k\to\infty$, $k\in K_1$, то найдется такой


номер $N_1$, что $x_k\in\omega\cap\cal D$ $\forall k\ge N_1$,


$k\in K_1$. Тогда


\begin{equation}


\begin{split}


\Bigl\la f'_0(x_k),\,\sum_{i\in I(\Bar x,0)}\Bar\alpha_i & f'_i(x_k)


\Bigr\ra+\frac{M}{2}\Bigl\|\,\sum_{i\in I(\Bar x,0)}


\Bar\alpha_if'_i(x_k)\Bigr\|^2\le \frac{\delta}{2},\\


f_j\Bigl(x_k & +\sum_{i\in I(\Bar x,0)}\Bar\alpha_if'_i(x_k)\Bigr)


\le 0 \quad \forall j\in J,


\end{split}


\end{equation}


$f'_i(x_k)\ne 0$ $\forall i\in I(\Bar x,0)$, $k\ge N_1$, $k\in K_1$.


По


лемме 3 существует такой номер $N_2\ge N_1$, что выполняется


включение (17) при всех $k\ge N_2$, $k\in K_1$. Положим


$\Bar\alpha_i=0$ $\forall i\in J\setminus J(\Bar x,0)$, тогда


$\sum\limits_{i\in I(\Bar x,0)}\Bar\alpha_if'_i(x_k)=


\sum\limits_{i\in I(x_k,\varepsilon_k)}\Bar\alpha_if'_i(x_k)$,


$k\ge N_2$, $k\in K_1$, и из (38) следуют неравенства


\begin{gather}


\Bigl\la f'_0(x_k),\,\sum_{i\in I(x_k,\varepsilon_k)}


\Bar\alpha_if'_i(x_k)


\Bigr\ra+\frac{M}{2}\Bigl\|\,\sum_{i\in I(x_k,\varepsilon_k)}


\Bar\alpha_if'_i(x_k)\Bigr\|^2\le \frac{\delta}{2},\\


f_j\Bigl(x_k+\sum_{i\in I(x_k,\varepsilon_k)}


\Bar\alpha_if'_i(x_k)\Bigr)\le 0


\quad \forall j\in J,


\end{gather}


$k\ge N_2$, $k\in K_1$.





Покажем далее, что случай, когда при некотором $k\in K$ одновременно


выполняются равенства


\begin{equation}


y_k=\argmin_{x\in\cal D_2(x_k,\varepsilon_k)}F_k(x),\quad


J(x_k,\varepsilon_k)=\varnothing,


\end{equation}


возможен лишь для конечного числа номеров $k\in K_1$.


Допустим противное, т.е.


существует такая подпоследовательность $K_2\subset K_1$, что (41)


справедливо для


всех $k\in K_2$. Тогда


$\cal D_2(x_k,\varepsilon_k)=M(x_k,\varepsilon_k)=


\{x\in R_n:x=x_k+\alpha_0f'_0(x_k),\;\alpha_0\in R_1\}$,


$y_k=x_k+\alpha_0(x_k)f'_0(x_k)$, где


$\alpha_0(x_k)=-\|f'_0(x_k)\|^2/\la f''_0(x_k)f'_0(x_k),\;


f'_0(x_k)\ra$, $k\in K_2$.


Значит, в силу (25)


$\alpha_0(x_k)\|f'_0(x_k)\|\to 0$, $k\to\infty$, $k\in K_2$, что


невозможно, поскольку справедливо (2) и


$0\le\la f''_0(x_k)f'_0(x_k),\;f'_0(x_k)\ra\le Md^2$


$\forall k\in K_2$ согласно (6),


(23), (16).


Таким образом, выбрать указанную подпоследовательность номеров $K_2$


нельзя, и найдется такой номер $N_3\ge N_2$, что для всех


$k\ge N_3$, $k\in K_1$, точки


$y_k$, а


значит, и числа $\alpha_i(x_k)$, $i\in I(x_k,\varepsilon_k)$,


получены с


помощью решения задачи (3),


(4) или (3), (5) при $\Bar x=x_k$, $\varepsilon=\varepsilon_k$.


Тогда при


всех $k\ge N_3$, $k\in K_1$, значения


$\alpha_i=\alpha_i(x_k)$, $i\in I(x_k,\varepsilon_k)$, удовлетворяют


ограничениям


\begin{equation}


f_j\Bigl(x_k+\sum_{i\in I(x_k,\varepsilon_k)}


\alpha_if'_i(x_k)\Bigr)\le 0


\end{equation}


для всех $j\in J$ или $j\in J(x_k,\varepsilon_k)$ в зависимости от


способа задания точки $y_k$. Согласно


(40) числа $\alpha_i=\Bar\alpha_i$, $i\in I(x_k,\varepsilon_k)$, также


удовлетворяют неравенствам (42) при $k\ge N_3$, $k\in K_1$.


Следовательно,


\begin{multline*}


\Bigl\la f'_0(x_k),\,\sum_{i\in I(x_k,\varepsilon_k)}\alpha_i(x_k)


f'_i(x_k)


\Bigr\ra+\frac{1}{2}\sum_{i,j\in I(x_k,\varepsilon_k)}


\alpha_i(x_k)\alpha_j(x_k)\la c_i(x_k),f'_j(x_k)\ra\le \\


\le \Bigl\la f'_0(x_k),\,\sum_{i\in I(x_k,\varepsilon_k)}


\Bar\alpha_if'_i(x_k)


\Bigr\ra+\frac{1}{2}\sum_{i,j\in I(x_k,\varepsilon_k)}


\Bar\alpha_i\Bar\alpha_j\la c_i(x_k),f'_j(x_k)\ra,


\quad k\ge N_3,\; k\in K_1.


\end{multline*}


Оценивая с помощью (6), (23) вторые слагаемые в обеих частях последнего


неравенства, получим


\begin{multline*}


\delta_k=


\Bigl\la f'_0(x_k),\,\sum_{i\in I(x_k,\varepsilon_k)}\alpha_i(x_k)


f'_i(x_k)


\Bigr\ra+\frac{\mu}{2}\Bigl\|\sum_{i\in I(x_k,\varepsilon_k)}


\alpha_i(x_k)f'_i(x_k)\Bigr\|^2 \le \\


\le \Bigl\la f'_0(x_k),\,\sum_{i\in I(x_k,\varepsilon_k)}


\Bar\alpha_if'_i(x_k)


\Bigr\ra+\frac{M}{2}\Bigl\|\sum_{i\in I(x_k,\varepsilon_k)}


\Bar\alpha_if'_i(x_k)\Bigr\|^2\quad


\forall k\ge N_3,\; k\in K_1.


\end{multline*}


Тогда из (39) будем иметь $\delta_k\le\delta/2$


$\forall k\ge N_3$, $k\in K_1$. Это неравенство


противоречит


предельным соотношениям (35). Таким образом, предположение о том, что


неравенство (24) выполняется как строгое, неверно. Значит,


последовательность $\{x_k\}$, $k\in K$, является минимизирующей, и из


условий


сильной выпуклости и непрерывности функции $f_0(x)$ на $\cal D$


следует,


что $x_k\to x^*$, $k\to\infty$.


\end{pf}





Пусть шаги в алгоритме выбираются из условия (14), и в дополнение к


условиям леммы 5 для $f''_0(x)$ на множестве $\cal D$


выполняется условие


Липшица с константой $L'$, т.е.


$$


\|f''_0(x)-f''_0(y)\|\le L'\|x-y\|\quad


\forall x,y\in \cal D.


$$


Покажем, что тогда найдется такой номер $N$, что


$\beta_k=1$ $\forall k\ge N$. Из (27), (28)


при $\beta=1$ с учетом условия Липшица и оценки (30) следуют неравенства


$f_0(x_k+s_k)-f_0(x_k)\le\tau_kF_k(y_k)+(L'/2)\|s_k\|^3\le


-\tau_k|F_k(y_k)|+(L'/\mu)\tau_k|F_k(y_k)|\,\|s_k\|$


$\forall k\in K$.


Согласно (25) $\|s_k\|\to 0$, $k\to\infty$, поэтому существует такой


номер $N$, что $(L'/\mu)\|s_k\|\le1-\eta$


$\forall k\ge N$. Значит,


$f_0(x_k)-f_0(x_k+s_k)\ge\tau_k|F_k(y_k)|(1-(L'/\mu)\|s_k\|)\ge


\eta\tau_k|F_k(y_k)|$


$\forall k\ge N$, т.е. условие (14) выполняется при $i=i_k=0$


$\forall k\ge N$,


что и требовалось доказать. Таким образом, если точка $y_k$ для всех


$k\in K$


выбирается первым способом, то $x_{k+1}=y_k$ $\forall k\ge N$.





Далее, отметим, что доказательство теоремы 2 не использует


непосредственно условия (14), (15) выбора шагов в алгоритме, а


опирается на результаты (24) и (25). Поэтому алгоритм допускает и


другие способы задания шагов $\beta_k$, если они гарантируют выполнение


условий (24), (25). Предполагая дополнительно, что


$f'_0(x)$ удовлетворяет на


множестве $\cal D$ условию Липшица с константой $L$, приведем ниже еще


два способа выбора чисел $\beta_k$.





Если известна константа $L$ или ее верхняя оценка, то в алгоритме можно


положить


$$


\beta_k=\min\{1,\rho_k\tau_k|F_k(y_k)|/\|s_k\|^2\},


$$


где


$0<\rho'\le\rho_k\le 2(1-\rho'')/L$,


$0<\rho''<1$, $\forall k\in K$. Покажем, что тогда


справедливы утверждения леммы 5. По


известному свойству выпуклых функций


$f_0(x_k)-f_0(x_{k+1})\ge-\beta_k\la f'_0(x_k),s_k\ra-


\beta^2_k L\|s_k\|^2/2$ $\forall k\in K$. Учитывая, что


$\la f'_0(x_k),s_k\ra=\tau_k\la f'_0(x_k),p_k\ra=


\tau_kF_k(y_k)-(\tau_k/2)\la f''_0(x_k)p_k,p_k\ra$, из


последнего неравенства и условия (6) получим оценку


$f_0(x_k)-f_0(x_{k+1})\ge\beta_k\tau_k|F_k(y_k)|-


\beta^2_k L\|s_k\|^2/2$ $\forall k\in K$. Отсюда при


$\beta_k=1\le\rho_k\tau_k|F_k(y_k)|/\|s_k\|^2$ следует неравенство


$$


f_0(x_k)-f_0(x_{k+1})\ge\rho''\tau_k|F_k(y_k)|,


$$


а при $\beta_k=\rho_k\tau_k|F_k(y_k)|/\|s_k\|^2<1$


--- неравенство


$$


f_0(x_k)-f_0(x_{k+1})\ge\rho'\rho''\tau^2_k|F_k(y_k)|^2/


\|s_k\|^2.


$$


Таким образом, последовательность $\{f_0(x_k)\}$ монотонно убывает, а


поскольку она


ограничена, то справедливы утверждение (24) и равенство (34). Как и в


лемме 2, при условии релаксационности процесса легко показать


ограниченность последовательностей $\{x_k\}$, $\{s_k\}$. Значит, из


последних двух


неравенств и (34) вытекает предельное соотношение


$\tau_k|F_k(y_k)|\to 0$, $k\to\infty$. Тогда из


оценки (30) и леммы 4 следует (25), т.е. теорема 2 справедлива при


данном способе задания шагов.





Пусть теперь шаги $\beta_k$ в алгоритме выбираются следующим образом


(\cite{Vas}, с.293):


\begin{gather*}


0\le\beta_k\le 1,\quad f_0(x_k+\beta_ks_k)=


\min_{0\le\beta\le 1}f_0(x_k+\beta s_k)+\sigma_k,\\


\sigma_k\ge 0,\quad \sum^\infty_{k=0}\sigma_k<\infty,\quad


\forall k\in K.


\end{gather*}





Заметим, что такой способ задания шагов не гарантирует релаксационность


процессу минимизации. Докажем при условии ограниченности множества


$\cal D$,


что утверждения (24), (25) справедливы и в этом случае. Действительно,


для всех $\beta\in[0,1]$, $k\in K$ с учетом (6) выполняются неравенства


$f_0(x_k)-f_0(x_{k+1})+\sigma_k\ge f_0(x_k)-f_0(x_k+\beta s_k)\ge


-\beta\la f'_0(x_k),s_k\ra-\beta^2 L\|s_k\|^2/2\ge\beta\tau_k


|F_k(y_k)|-\beta^2 L\|s_k\|^2/2$.


Так как $x_k\in\cal D$,


$k\in K$, то последовательности $\{x_k\}$, $\{s_k\}$ ограничены, и,


значит, $\|s_k\|^2\le d'<\infty$ $\forall k\in K$. Тогда из


последних неравенств и оценки (30) следует


\begin{equation}


\|s_k\|^2\le(f_0(x_k)-f_0(x_{k+1})+\sigma_k)2/\beta\mu+


\beta Ld'/\mu


\end{equation}


$\forall\beta\in(0,1]$, $k\in K$. Поскольку


$f_0(x_{k+1})\le f_0(x_k)+\sigma_k$, $k\in K$, и


последовательность $\{f_0(x_k)\}$ ограничена снизу, то


существует (\cite{Vas}, с.93) конечный предел (24), а значит,


выполняется (34). Тогда, в силу $\sigma_k\to 0$, $k\to\infty$, из (43)


следуют неравенства


$0\le\underset{k\to\infty}{\varliminf}{}\|s_k\|^2\le


\underset{k\to\infty}{\varlimsup}{}\|s_k\|^2\le\beta Ld'/\mu$


при каждом фиксированном $\beta\in(0,1]$. Устремляя


$\beta\to 0+$, получим из этих


неравенств


утверждение (25). Таким образом, если $\cal D$ ограничено,


то алгоритм с указанным способом выбора шагов является сходящимся.





Отметим, что если $f_j(x)=\la a_j,x\ra$, $a_j\in R_n$, $j\in J$, то


задача отыскания точки


$y_k$ на шаге 2 алгоритма может быть решена за конечное число итераций


(см., напр, \cite{Ps}, \cite{Su}).





При $J(x_k,\varepsilon_k)=\varnothing$ направление спуска на


$k$-й итерации


алгоритма выгоднее получать


с помощью задачи (3), (5) при


$\Bar x=x_k$, $\varepsilon=\varepsilon_k$. Ее решение в указанном


случае не


представляет труда, поскольку


$I(x_k,\varepsilon_k)=\{0\}$,


$\cal D_2(x_k,\varepsilon_k)=\{x\in R_n:x=x_k+\alpha_0f'_0(x_k),\;


\alpha_0\in R_1\}$, и


значение $\alpha_0(x_k)$ находится по явной


формуле. В методе Ньютона, например, в такой же ситуации для построения


направления приходится решать вспомогательную задачу условной


минимизации. Кроме того, еще раз подчеркнем, что число переменных и


число ограничений вспомогательных задач выбора направления в


предложенном алгоритме зависят лишь от количества активных на каждом


шаге ограничений исходной задачи. Поэтому при $m<n$ алгоритм имеет


некоторые преимущества перед известными методами второго порядка


(см., напр, \cite{Vas}, \cite{Ps}, \cite{Gill}), в которых на каждой


итерации аналогичная задача выбора подходящего направления имеет $n$


переменных и $m$ ограничений.





\begin{thebibliography}{9}





\bibitem{Vas}


Васильев Ф.П. {\em Численные методы решения экстремальных задач}. --


М.: Наука, 1988. -- 552 с.


\bibitem{Zab}


Заботин И.Я. {\em Метод условной минимизации с параметрическим заданием


подходящих направлений} // Изв. вузов. Математика. -- 1996. -- \No\,12.


-- С.17--29.


\bibitem{Kar}


Карманов В.Г. {\em Математическое программирование}. -- М.: Наука,


1986. -- 288 с.


\bibitem{Ps}


Пшеничный Б.Н., Данилин Ю.М. {\em Численные методы в экстремальных


задачах}. -- М.: Наука, 1975. -- 320 с.


\bibitem{Su}


Сухарев А.Г., Тимохов А.В., Федоров В.В. {\em Курс методов


оптимизации}. -- М.: Наука, 1986. -- 328 с.


\bibitem{Gill}


{\em Численные методы условной оптимизации}. Ред. Гилл Ф., Мюррей У. --


М.: Мир, 1977. -- 290~с.





\end{thebibliography}





\vskip 2cm





\post{\it Казанский государственный университет}{\it Поступила}


\post{}{23.05.1997}





\end{document}


